\unitopen{41}
\begin{candos}
\item budować złożone argumenty spójnikami i przysłówkami łączącymi poziomu B2;
\item używać spójników dwuczęściowych (\de{je \dots{} desto}, \de{nicht nur \dots{} sondern auch}, \de{weder \dots{} noch});
\item zamieniać zdanie podrzędne na konstrukcję z \de{ohne zu} i \de{statt zu};
\item napisać ustrukturyzowany tekst argumentacyjny z wnioskiem.
\end{candos}

\section{Słownictwo}
\vocabtable{41}

\section{Cztery rodziny łączników}
\begin{center}\small
\begin{tabular}{@{}p{0.18\linewidth}p{0.36\linewidth}p{0.36\linewidth}@{}}
\toprule
\textsf{\textbf{Funkcja}} & \textsf{\textbf{Zdanie podrzędne (czasownik na końcu)}} & \textsf{\textbf{Przysłówek (czasownik 2.)}}\\
\midrule
sposób & \de{indem man den Cache nutzt} & ---\\
skutek & \de{sodass die Latenz sinkt} & \de{folglich, somit, infolgedessen}\\
warunek & \de{falls}, \de{sofern}, \de{es sei denn, \dots} & \de{andernfalls, sonst}\\
przeciwstawienie & \de{während}, \de{wohingegen} & \de{hingegen, dagegen, dennoch}\\
\bottomrule
\end{tabular}
\end{center}
Przykłady: \de{Wir haben die Latenz gesenkt, \textbf{indem} wir einen Cache eingeführt haben.} \de{Die Abfrage ist optimiert, \textbf{sodass} sie
nur noch zwei Sekunden dauert.} \de{Das Release erfolgt am Freitag, \textbf{es sei denn}, die Tests schlagen fehl.} \de{Die Kosten sind gestiegen, \textbf{dennoch}
lohnt sich die Migration.}

\section{Spójniki dwuczęściowe}
\de{je \dots{} desto/umso}: \de{Je mehr Daten wir haben, desto besser wird das Modell.} \de{nicht nur \dots{} sondern auch}: \de{Der Dienst ist nicht nur
schneller, sondern auch günstiger.} \de{weder \dots{} noch}: \de{Er hat weder Zeit noch Lust.} \de{entweder \dots{} oder}: \de{Entweder
wir verschieben, oder wir arbeiten am Wochenende.} \de{sowohl \dots{} als auch}: \de{Wir brauchen sowohl Tests als auch Doku.}
\de{zwar \dots{} aber}: \de{Das ist zwar teuer, aber es lohnt sich.}

\section{Zamiana na bezokolicznik}
Gdy podmiot jest ten sam: \de{ohne dass er fragt $\rightarrow$ ohne zu fragen}; \de{anstatt dass wir warten $\rightarrow$ statt zu warten}.

\begin{falle}
Po polsku „im więcej, tym lepiej” nie wymaga zmiany szyku, a po niemiecku \de{je} otwiera zdanie podrzędne (czasownik na końcu), a
\de{desto} główne (odwrócony szyk: \de{desto besser wird das Modell}). Błąd: \de{je mehr wir haben, desto wir sind besser}. Drugi:
\de{falls} i \de{wenn} nie są zamienne: \de{falls} to ewentualność, \de{wenn} to także czas. Trzeci: pomylenie \de{dennoch}
(bez zmiany szyku zdania podrzędnego) z \de{obwohl}.
\end{falle}

\begin{szwajcaria}
W szwajcarskim piśmie biznesowym łączniki \de{sofern}, \de{falls}, \de{sodass} i \de{somit} są częste; \de{somit} otwiera wnioski.
W dyskusjach ustnych natomiast Szwajcarzy używają prostszych: \de{also}, \de{dann}, \de{aber}. Pisząc raport, wybierz przysłówki z
tabeli; mówiąc, trzymaj się prostoty.
\end{szwajcaria}

\begin{praca}
\de{Wir haben die Ladezeit verkürzt, indem wir die Bilder komprimiert haben, sodass die Seite nun doppelt so schnell lädt.}
\de{Je früher wir testen, desto weniger Fehler gelangen in Produktion.} \de{Das Projekt ist nicht nur technisch, sondern auch organisatorisch
anspruchsvoll.}
\end{praca}

\section{Dialogi}
\textbf{Dialog 1 --- argumentacja}
\begin{dialog}
\turn{A}{Wir sollten automatisch testen, anstatt jeden Fall manuell zu prüfen.}{Powinniśmy testować automatycznie, zamiast sprawdzać każdy przypadek ręcznie.}
\turn{B}{Das stimmt zwar, aber der Aufwand ist hoch. Es sei denn, wir nutzen bestehende Tools.}{To prawda, ale nakład jest duży. Chyba że użyjemy istniejących narzędzi.}
\turn{A}{Je früher wir anfangen, desto weniger kostet es uns später.}{Im wcześniej zaczniemy, tym mniej nas to później kosztuje.}
\turn{B}{Dann schlage ich vor, dass wir einen Pilot starten, sodass wir den Nutzen messen können.}{Wobec tego proponuję start pilota, żebyśmy mogli zmierzyć korzyść.}
\end{dialog}

\section{Ćwiczenia}
\begin{cwiczenia}
\exitem{Połącz (\de{indem}): \de{Wir senken die Latenz. Wir führen einen Cache ein.}}{\de{Wir senken die Latenz, indem wir einen Cache einführen.}}
\exitem{Połącz (\de{sodass}): \de{Die Abfrage ist optimiert. Sie dauert zwei Sekunden.}}{\de{Die Abfrage ist optimiert, sodass sie nur zwei Sekunden dauert.}}
\exitem{Połącz (\de{je \dots{} desto}): \de{viele Daten / ein gutes Modell}.}{\de{Je mehr Daten wir haben, desto besser wird das Modell.}}
\exitem{Zamiast zdania podrzędnego: \de{ohne dass er fragt}.}{\de{ohne zu fragen}}
\exitem{Zamiast zdania podrzędnego: \de{anstatt dass wir warten}.}{\de{statt zu warten}}
\exitem{Użyj \de{nicht nur \dots{} sondern auch}: \de{schnell / günstig}.}{\de{Der Dienst ist nicht nur schneller, sondern auch günstiger.}}
\exitem{Użyj \de{weder \dots{} noch}: \de{Zeit / Lust}.}{\de{Er hat weder Zeit noch Lust.}}
\exitem{Popraw: \de{Je mehr wir haben, desto wir sind besser.}}{\de{Je mehr wir haben, desto besser sind wir.}}
\exitem{Wskaż: \de{dennoch} czy \de{obwohl} wymaga czasownika na końcu?}{\de{obwohl} (zdanie podrzędne); po \de{dennoch} czasownik jest na 2.\ miejscu}
\exitem{Przetłumacz: „Test powiedzie się, chyba że baza nie odpowie.”}{\de{Der Test ist erfolgreich, es sei denn, die Datenbank antwortet nicht.}}
\end{cwiczenia}

\weekend{Napisz tekst argumentacyjny (250 słów) „Monolit czy mikroserwisy?”, używając co najmniej sześciu łączników z tej jednostki
i wniosku z \de{somit} lub \de{folglich}. Powiedz jego streszczenie w dwóch minutach.}
